\chapter*{卷三\quad 撤离}
\phantomsection
\addcontentsline{toc}{chapter}{卷三\quad 撤离}
\markboth{卷三\quad 撤离}{卷三\quad 撤离}

\noindent
卷三十二片，是全书最重要的一卷。它们记的是公元前二百二十一前的最后一次清点。

《战争的雷霆》第一章所引「第 2187 号卷宗第 121 叶，设施退役清单，共十九项」，所指即简二十一。该书只引了第七项，此处给出全文。

---

\begin{jian}{简二十}
元年，会计。在籍十一万九千。可去者十一万九千。留者无。
\end{jian}
\chuchu{原件残片 57 号，长 22.0 cm，金属。}
\xiaojì{「元年」何指，不可考。若指某纪年之始，则此片当在他片之前。「留者无」三字刻痕较深，与他字不类，或为后加。}

\begin{jian}{简二十一}
退役清单，凡十九项。第一项：居住环六，封。第二项：气路，断。第三项：水路，断。第四项：库，封。第五项：籍，焚。第六项：器，毁。第七项：垂直运输装置（电梯）已拆，支柱于四十米高度截断，截断处标记，勿挖。第八项：界，断。第九项：外舍，平。第十项：垣门，塞。
\end{jian}
\chuchu{原件残片 60 号，长 28.4 cm，为全书最长一片，下端接续处断裂。}
\xiaojì{全文十九项，此片存十项，其余九项据残片 61 号补，见简二十二。「勿挖」二字为当时俗语，谓此处所埋之物不可开掘，非警告后人之辞。}

\begin{jian}{简二十二}
第十一项：山，封。第十二项：炉，熄。第十三项：药，弃。第十四项：册，收。第十五项：界石，移。第十六项：标，除。第十七项：名，削。第十八项：口，减。第十九项：门，闭。
\end{jian}
\chuchu{原件残片 61 号，长 23.7 cm，金属，字迹较前片为浅。}
\xiaojì{「口，减」，谓登记人口数予以削减，非谓减少人口。「名，削」与之相接，谓削去名册。「药，弃」之药为何，全书未再出现，存疑。}

\begin{jian}{简二十三}
支柱截，深四十米。填以土，夯实。上覆石板三层，石板之上植草，草不活。
\end{jian}
\chuchu{原件残片 63 号，长 16.4 cm，陶瓷。}
\xiaojì{「草不活」，三字独立成句，为他片所未见。所指为其地不宜草木，抑或为当时实况，不可辨。}

\begin{jian}{简二十四}
装载：人十一万九千。粟，尽。水，尽。气，尽。器，择要者三千件。册，尽。
\end{jian}
\chuchu{原件残片 65 号，长 20.1 cm，金属。}
\xiaojì{「人十一万九千」，与简二十合。「粟尽、水尽、气尽」三语，谓库存全部带走，不留后用。此与简二十一「第十二项：炉，熄」相应，可见撤离之后此城不再维持。}

\begin{jian}{简二十五}
凡去者，各给符一枚。符不还。符上刻其名。无名者不得登。
\end{jian}
\chuchu{原件残片 67 号，长 14.2 cm，陶瓷。}
\xiaojì{「符不还」，谓符为个人之物，离城后仍归其所有。「无名者不得登」，与前片「名，削」矛盾：既削其名，何由得登。此为本卷最大的一处不可解，详下编第六章。}

\begin{jian}{简二十六}
世代存活率：初代六一四，二代三八七，三代一二九，四代四二。四代之后，无记。
\end{jian}
\chuchu{原件残片 69 号，长 18.7 cm，金属，表面有铅笔痕一道，非刻。}
\xiaojì{四个数字以千分率计，即 61.4\%、38.7\%、12.9\%、4.2\%，与《战争的雷霆》第一章第四节所引完全一致。「四代之后无记」，为全书最著名的空白。}

\begin{jian}{简二十七}
地热之衰，二百年间，百三十一分。始缓，终急。
\end{jian}
\chuchu{原件残片 71 号，长 15.3 cm，金属。}
\xiaojì{「百三十一分」即 3.1\%，与《战争的雷霆》所称二百年间下降 3.1\% 合。此片为撤离近因之唯一原件依据。}

\begin{jian}{简二十八}
外工三百二十人，不入装载之数。留。
\end{jian}
\chuchu{原件残片 73 号，长 10.8 cm，金属，下端蚀尽。}
\xiaojì{「留」字后有缺文，所缺或为一句，或为一字。若此三百二十人确被留下，则简二十「留者无」为误记。此点无法判定。}

\begin{jian}{简二十九}
六批之数：一批六千，二批八千，三批一万，四批一万三千，五批一万八千，六批六万五千。凡十二万。
\end{jian}
\chuchu{原件残片 75 号，长 25.6 cm，金属，字口深，为他片之最。}
\xiaojì{六批之数为本书最完整的一组数据，与《战争的雷霆》第一章所复原者合。「凡十二万」，与简二十、简二十四「十一万九千」相差一千，或为取整，或为另计。}

\begin{jian}{简三十}
六批者，非一时之事。前五批，岁计之内。第六批，岁计之外。
\end{jian}
\chuchu{原件残片 77 号，长 17.9 cm，陶瓷。}
\xiaojì{「岁计之内」谓前五批按常规计划执行；「岁计之外」谓第六批不在原计划内。此片直接支持《战争的雷霆》所说「前 221 年是最后一批，也是唯一一批走得仓促的」。}

\begin{jian}{简三十一}
门闭之日，无送者，无哭声。闭者七人，回首一次。
\end{jian}
\chuchu{原件残片 79 号，长 12.9 cm，金属，为卷三最后一片。}
\xiaojì{「回首一次」四字为全书唯一一处带动作的记载，与他片之纯记数者不类。此片是否与撤离诸片同批，存疑。}
